\documentclass[12pt,a4paper]{article}

% Paquetes
\usepackage[utf8]{inputenc}
\usepackage[T1]{fontenc}
\usepackage[spanish,es-nodecimaldot]{babel}
\usepackage[a4paper,margin=2.5cm]{geometry}
\usepackage{amsmath}
\usepackage{amssymb}
\usepackage{booktabs}
\usepackage{float}
\usepackage{graphicx}
\usepackage{siunitx}
\usepackage{tikz}
\usepackage{pgfplots}
\usepackage[hidelinks]{hyperref}
\usepackage{url}

\usetikzlibrary{arrows.meta,positioning,shapes.geometric}
\pgfplotsset{compat=1.18}

\graphicspath{{../figuras/}}

\sisetup{
    output-decimal-marker = {,},
    group-separator = {.},
    group-minimum-digits = 4
}

% Bibliografía
\usepackage[
    backend=biber,
    style=ieee,
    sorting=none
]{biblatex}

\addbibresource{references.bib}

\begin{document}

% Portada
\begin{titlepage}
    \centering

    {\Large\textbf{UNIVERSIDAD TÉCNICA ESTATAL DE QUEVEDO}\par}
    \vspace{0.4cm}

    {\large Facultad de Ciencias de la Computación\par}
    {\large Carrera de Ingeniería de Software\par}

    \vfill

    {\Large\textbf{TA-IND-04}\par}
    \vspace{0.4cm}

    {\LARGE\textbf{Análisis de Rendimiento Paralelo aplicado al Proyecto Fin de Curso}\par}

    \vfill

    \begin{flushleft}
        \textbf{Asignatura:} Aplicaciones Distribuidas (ISR-701)\\[0.2cm]
        \textbf{Unidad:} Unidad 4 -- Cómputo Paralelo y Distribuido\\[0.2cm]
        \textbf{Estudiante:} Jeremy Ruperto Gaibor Rodríguez\\[0.2cm]
        \textbf{Equipo PE-U4:} Equipo C\\[0.2cm]
        \textbf{PFC de referencia:} Tienda Tech\\[0.2cm]
        \textbf{Transformación foco:} T1 -- Filtrado y selección\\[0.2cm]
        \textbf{Período académico:} 2026--2027 PPA\\[0.2cm]
        \textbf{Docente:} Gleiston C. Guerrero-Ulloa, M.Sc.\\[0.2cm]
        \textbf{Repositorio individual:}\\
        \url{AQUI_VA_LA_URL_DE_TU_REPOSITORIO}\\[0.2cm]
        \textbf{Fecha:} Agosto de 2026
    \end{flushleft}

    \vfill
\end{titlepage}

% Introducción
\section{Introducción y contexto}

El procesamiento paralelo y distribuido permite repartir una carga de trabajo
entre varias unidades de cómputo. Sin embargo, aumentar la cantidad de recursos
no garantiza una reducción proporcional del tiempo de ejecución. En Apache
Spark, la cantidad de executors constituye un parámetro relevante y una
asignación inadecuada puede producir un uso poco eficiente de los recursos o
un mayor tiempo de ejecución \cite{sen2021}. Asimismo, el ajuste de los
parámetros de configuración de Spark constituye un problema de optimización
que depende de las características de cada aplicación \cite{cheng2021}.

El presente informe analiza de manera individual los resultados obtenidos en la
práctica PE-U4. En el experimento se implementaron cinco transformaciones
equivalentes mediante pandas y PySpark sobre un conjunto sintético de
\num{500000} solicitudes y 14 columnas. Para cada configuración se realizaron
cinco repeticiones y se utilizó la mediana como tiempo representativo. Antes de
las repeticiones registradas se descartó una ejecución inicial de calentamiento.

La transformación seleccionada como foco individual es T1, correspondiente al
filtrado y selección. Esta operación procesa los \num{500000} registros de
entrada y produce \num{33117} registros como resultado.

Los datos utilizados son trazables al repositorio PE-U4 del Equipo C:

\begin{center}
\url{https://github.com/ffarinangog2/pe-u4-spark-equipo-c}
\end{center}

\noindent
Las mediciones adicionales de T1 con uno y dos executors se incorporaron en el
commit \texttt{6b363ab}. Los resultados se encuentran registrados en
\texttt{resultados/tiempos\_crudos.csv} y
\texttt{resultados/tiempos\_resumen.csv}.

El entorno documentado en el proyecto utiliza PySpark 3.5.0, Python 3.12.10,
OpenJDK 17.0.20 y Windows 11. Spark fue ejecutado en modo Standalone con el
master \texttt{spark://127.0.0.1:7077}. Cada executor fue configurado con un
core y \SI{512}{\mebi\byte} de memoria. Además, el número de particiones, el
paralelismo por defecto y el máximo de cores se ajustaron al número de
executors utilizado en cada configuración.

% Análisis de speedup
\section{Análisis propio de los resultados de speedup}
\label{sec:speedup}

La comparación general entre pandas y PySpark utiliza el tiempo mediano de
pandas y el tiempo obtenido por Spark con cuatro executors. El speedup se
determina mediante:

\[
S=
\frac{T_{\text{pandas}}}
{T_{\text{spark}}}.
\]

La Tabla~\ref{tab:general} resume los resultados de las cinco transformaciones.

\begin{table}[H]
\centering
\caption{Resultados generales de PE-U4 con cuatro executors.}
\label{tab:general}
\begin{tabular}{lrrrr}
\toprule
\textbf{Transf.} &
\textbf{$T_{\text{pandas}}$ (s)} &
\textbf{$T_{\text{spark}}$ (s)} &
\textbf{$S$} &
\textbf{$E$} \\
\midrule
T1 & 0.2119435 & 0.3456639 & 0.6131 & 0.1533 \\
T2 & 0.1171350 & 0.6874908 & 0.1704 & 0.0426 \\
T3 & 0.2888583 & 0.8657363 & 0.3337 & 0.0834 \\
T4 & 2.2658870 & 1.8100053 & 1.2519 & 0.3130 \\
T5 & 1.5252129 & 1.3564848 & 1.1244 & 0.2811 \\
\bottomrule
\end{tabular}
\end{table}

Los resultados muestran que T4 y T5 alcanzaron un speedup superior a uno,
mientras que T1, T2 y T3 fueron más rápidas mediante pandas. En el caso de T1
se obtuvo un speedup de \num{0.6131}, por lo que la ejecución de Spark con
cuatro executors necesitó más tiempo que pandas.

Este resultado no es inesperado en un motor distribuido. Sen et al. muestran
que distintas consultas de Spark pueden requerir niveles diferentes de
paralelismo y que incrementar la cantidad de executors no garantiza siempre
una reducción del tiempo de ejecución \cite{sen2021}. Por su parte, Cheng
et al. estudian el ajuste de parámetros de Spark como un problema de
optimización dependiente de la carga de trabajo \cite{cheng2021}.

\subsection{Escalado de T1}

Para estudiar el escalado de T1 se toma como referencia la ejecución Spark con
un executor:

\[
T_1=\SI{0.3134846}{\second}.
\]

Los resultados obtenidos fueron los siguientes:

\begin{table}[H]
\centering
\caption{Escalado experimental de T1.}
\label{tab:escalado}
\begin{tabular}{rrrr}
\toprule
\textbf{$N$} &
\textbf{$T_N$ (s)} &
\textbf{$S(N)$} &
\textbf{$E(N)$} \\
\midrule
1 & 0.3134846 & 1.0000 & 1.0000 \\
2 & 0.8624843 & 0.3635 & 0.1817 \\
4 & 0.3456639 & 0.9069 & 0.2267 \\
\bottomrule
\end{tabular}
\end{table}

Para dos executors, el speedup es:

\begin{align}
S(2)
&=
\frac{T_1}{T_2} \\
&=
\frac{0.3134846}{0.8624843} \\
&=
0.3635.
\end{align}

Para cuatro executors:

\begin{align}
S(4)
&=
\frac{T_1}{T_4} \\
&=
\frac{0.3134846}{0.3456639} \\
&=
0.9069.
\end{align}

En ninguna de las configuraciones se obtuvo una aceleración respecto de la
ejecución con un solo executor. La configuración con dos executors presentó el
peor resultado.

Sus cinco tiempos fueron \SI{0.8624843}{\second},
\SI{0.3758325}{\second}, \SI{1.1714799}{\second},
\SI{0.2754539}{\second} y \SI{1.4906592}{\second}. La diferencia entre el
menor y el mayor valor evidencia una variabilidad considerable en esta
configuración.

\subsection{Contraste con la Ley de Amdahl}

La Ley de Amdahl establece el límite del speedup que puede alcanzarse cuando
una fracción $p$ de un programa puede ejecutarse en paralelo
\cite{amdahl1967}:

\begin{equation}
S(N)=
\frac{1}
{(1-p)+\dfrac{p}{N}},
\qquad
S_{\max}=
\frac{1}{1-p}.
\label{eq:amdahl}
\end{equation}

El experimento no midió de forma independiente la fracción paralelizable $p$.
Por ello, se utiliza $p=1$ únicamente como una referencia ideal superior, es
decir, como el caso en el que no existe fracción serial ni sobrecarga. Esta
elección no implica afirmar que T1 sea completamente paralelizable durante una
ejecución real.

Bajo esa referencia ideal:

\[
S_{\text{teo}}(N)=N.
\]

Para dos executors, la desviación relativa entre el valor medido y la
referencia ideal es:

\begin{align}
\Delta_2
&=
\frac{S_{\text{med}}-S_{\text{teo}}}
{S_{\text{teo}}}
\times100 \\
&=
\frac{0.3635-2}{2}\times100 \\
&=
-81.83\%.
\end{align}

Para cuatro executors:

\begin{align}
\Delta_4
&=
\frac{0.9069-4}{4}\times100 \\
&=
-77.33\%.
\end{align}

Por tanto, los resultados experimentales se encuentran claramente por debajo
del límite ideal de Amdahl. El speedup real es incluso menor que uno para las
dos configuraciones.

Esto no contradice la Ley de Amdahl. La expresión representa un límite teórico
y no incorpora la sobrecarga introducida por un motor distribuido real. En
Spark existe trabajo adicional relacionado con la creación, planificación y
coordinación de las tareas.

\begin{figure}[H]
\centering

\begin{tikzpicture}
\begin{axis}[
    width=0.82\textwidth,
    height=7.5cm,
    xlabel={Número de executors $N$},
    ylabel={Speedup $S(N)$},
    xtick={1,2,4},
    grid=major,
    legend pos=north west
]

\addplot[
    mark=*,
    thick
]
coordinates {
    (1,1)
    (2,2)
    (4,4)
};
\addlegendentry{Amdahl ideal ($p=1$)}

\addplot[
    mark=square*,
    thick
]
coordinates {
    (1,1)
    (2,0.3635)
    (4,0.9069)
};
\addlegendentry{T1 experimental}

\end{axis}
\end{tikzpicture}

\caption{Speedup experimental de T1 frente al límite ideal de Amdahl. Elaboración propia.}
\label{fig:speedup}
\end{figure}

La Fig.~\ref{fig:speedup} muestra que el comportamiento experimental no sigue
el incremento esperado en un escenario ideal. Con dos executors el rendimiento
cae de forma considerable. Al utilizar cuatro existe una recuperación, pero el
tiempo sigue siendo superior al registrado con un solo executor.

% Diagnóstico
\section{Diagnóstico de la fracción no escalable}

La métrica de Karp-Flatt permite estimar experimentalmente una fracción serial
a partir del speedup observado y del número de unidades de procesamiento
\cite{karpflatt1990}:

\begin{equation}
e=
\frac{
\dfrac{1}{S}
-
\dfrac{1}{N}
}
{
1-\dfrac{1}{N}
},
\qquad N>1.
\label{eq:karp}
\end{equation}

Para dos executors:

\begin{align}
e_2
&=
\frac{
\dfrac{1}{0.3635}
-
\dfrac{1}{2}
}
{
1-\dfrac{1}{2}
} \\
&\approx
4.5026.
\end{align}

Para cuatro executors:

\begin{align}
e_4
&=
\frac{
\dfrac{1}{0.9069}
-
\dfrac{1}{4}
}
{
1-\dfrac{1}{4}
} \\
&\approx
1.1369.
\end{align}

Los dos valores son superiores a uno. En estas condiciones no pueden
interpretarse como una fracción serial física convencional, ya que el modelo
está reflejando un escenario en el que el tiempo paralelo es peor que el tiempo
de referencia.

Además, $e$ disminuye de \num{4.5026} con dos executors a \num{1.1369} con
cuatro. Por tanto, la pérdida observada no se comporta como una fracción serial
constante. El resultado con dos executors está especialmente afectado por la
variabilidad y la sobrecarga de la ejecución.

\subsection{Sobrecarga observada}

El primer mecanismo que puede relacionarse con los resultados es la
planificación y coordinación del trabajo. Spark necesita preparar y distribuir
las tareas antes de ejecutarlas, por lo que el nivel de paralelismo elegido
puede afectar directamente al rendimiento. AutoExecutor estudia precisamente
cómo distintas consultas Spark SQL requieren diferente cantidad de executors y
cómo una selección inadecuada puede perjudicar el rendimiento
\cite{sen2021}.

El segundo mecanismo corresponde a la configuración de recursos. En el
experimento, \texttt{spark.executor.instances},
\texttt{spark.cores.max},
\texttt{spark.sql.shuffle.partitions} y
\texttt{spark.default.parallelism} cambian de acuerdo con $N$. Cheng et al.
muestran que la configuración de Spark tiene una influencia suficiente sobre
el rendimiento como para justificar métodos específicos de optimización de
parámetros \cite{cheng2021}.

El tercer elemento es la variabilidad observada en las propias mediciones. Con
dos executors se registraron tiempos comprendidos entre aproximadamente
\SI{0.275}{\second} y \SI{1.491}{\second}. Esta dispersión indica que la
sobrecarga efectiva de la configuración no se mantuvo constante durante las
cinco repeticiones.

La eficiencia del paralelismo se determina mediante:

\begin{equation}
E(N)=\frac{S(N)}{N}.
\label{eq:eficiencia}
\end{equation}

Para dos executors:

\[
E(2)=
\frac{0.3635}{2}
=
0.1817.
\]

Para cuatro:

\[
E(4)=
\frac{0.9069}{4}
=
0.2267.
\]

Por tanto, la eficiencia fue aproximadamente
\SI{18.17}{\percent} con dos executors y
\SI{22.67}{\percent} con cuatro. La segunda configuración presenta una mejora
respecto a la primera, pero ambas permanecen lejos de un aprovechamiento ideal
del paralelismo.

% Integración al PFC
\section{Propuesta de integración al PFC}

En el PFC Tienda Tech, el procesamiento distribuido se propone para el análisis
histórico de pedidos, ventas y movimientos de inventario. No se plantea su uso
dentro del flujo transaccional encargado de registrar directamente una compra,
porque dicho proceso requiere respuestas rápidas y trabaja con una cantidad
reducida de información por solicitud.

La integración se realizaría mediante procesamiento por lotes. Periódicamente
se extraerían de PostgreSQL los registros históricos necesarios y se
entregarían a un proceso Spark. Este proceso podría agrupar las ventas por
producto y período, analizar el movimiento del inventario y generar indicadores
de demanda.

La salida sería almacenada como información analítica preparada para consulta.
El módulo de reportes utilizaría estos resultados para evitar el procesamiento
completo del histórico cada vez que un usuario solicite información.

La decisión funcional soportada por esta integración sería la priorización de
la reposición del inventario. A partir de los resultados podrían identificarse
productos con alta demanda, bajo movimiento o necesidad de reposición.

Se propone una ejecución por lotes porque este procesamiento no necesita
realizarse durante la compra de un cliente. Su ejecución podría realizarse
periódicamente cuando se haya acumulado nueva información suficiente.

\begin{figure}[H]
\centering

\begin{tikzpicture}[
    node distance=1.15cm,
    every node/.style={
        draw,
        rounded corners,
        align=center,
        minimum width=3.4cm,
        minimum height=0.9cm
    },
    >={Latex}
]

\node (micro) {Microservicios\\Tienda Tech};
\node (postgres) [below=of micro] {PostgreSQL\\Datos operacionales};
\node (extract) [below=of postgres] {Extracción\\por lotes};
\node (spark) [below=of extract] {Apache Spark\\Procesamiento analítico};
\node (result) [below=of spark] {Resultados\\analíticos};
\node (report) [below=of result] {Reportes y apoyo\\a reposición};

\draw[->] (micro) -- (postgres);
\draw[->] (postgres) -- (extract);
\draw[->] (extract) -- (spark);
\draw[->] (spark) -- (result);
\draw[->] (result) -- (report);

\end{tikzpicture}

\caption{Propuesta de integración del procesamiento distribuido en Tienda Tech. Elaboración propia.}
\label{fig:arquitectura}
\end{figure}

Como se observa en la Fig.~\ref{fig:arquitectura}, Spark se ubica después de
la persistencia de los datos operacionales. De esta forma, el análisis se
mantiene separado de los microservicios responsables de las operaciones
transaccionales.

% Dimensionamiento
\section{Dimensionamiento y viabilidad}

El modelo propuesto para estimar la rentabilidad del procesamiento distribuido
considera un tiempo secuencial:

\[
T_{\text{sec}}(n)=\alpha n,
\]

y un tiempo distribuido:

\[
T_{\text{dist}}(n)=
\beta+
\frac{\gamma n}{N}.
\]

El umbral se expresa mediante:

\begin{equation}
n^{*}=
\frac{\beta}
{\alpha-\dfrac{\gamma}{N}},
\qquad
\alpha>
\frac{\gamma}{N}.
\label{eq:umbral}
\end{equation}

Los parámetros $\alpha$, $\beta$ y $\gamma$ deben obtenerse mediante ajuste
sobre mediciones realizadas con diferentes volúmenes de entrada. Los datos
actuales de T1 corresponden a un único volumen de \num{500000} registros, por
lo que no permiten identificar de manera independiente los tres parámetros del
modelo.

Por esta razón, con la evidencia disponible no se presenta un valor numérico
de $n^{*}$. Calcularlo utilizando únicamente las configuraciones de uno, dos y
cuatro executors implicaría utilizar variaciones de $N$ como si fueran
variaciones del volumen $n$, lo cual no corresponde al modelo de la
Ec.~\eqref{eq:umbral}.

Para completar el dimensionamiento sería necesario ejecutar T1 con distintos
tamaños de entrada manteniendo controladas las demás condiciones
experimentales. Gulino et al. estudian la predicción del tiempo de ejecución
de aplicaciones Spark considerando, entre otros elementos, el tamaño de los
datos y los recursos empleados \cite{gulino2020}. De forma relacionada,
Lattuada et al. analizan la asignación de recursos de aplicaciones Spark
mediante modelos de rendimiento orientados a determinar la capacidad necesaria
para alcanzar determinados objetivos de ejecución \cite{lattuada2020}.

Mientras no se determine experimentalmente el umbral, para volúmenes similares
a los evaluados se mantendría una solución basada en pandas o en consultas SQL,
ya que T1 no presentó una ventaja de rendimiento mediante Spark.

Para una futura implementación, el procesamiento distribuido se ejecutaría como
una tarea independiente por lotes y no como un servicio que deba permanecer
activo de forma permanente. Entre los principales riesgos se encuentran la
variabilidad del rendimiento, el costo de mantener infraestructura adicional,
un crecimiento de datos insuficiente para justificar Spark y la sincronización
de información procedente de varios microservicios.

% Postura crítica
\section{Postura crítica y alternativas}

Las mediciones disponibles indican que Spark no es actualmente la opción más
eficiente para T1 con \num{500000} registros. pandas obtuvo una mediana de
\SI{0.2119435}{\second}, mientras que Spark con cuatro executors necesitó
\SI{0.3456639}{\second}. Esto produjo un speedup de \num{0.6131}.

La evidencia experimental también muestra que aumentar la cantidad de
executors no produjo una mejora progresiva. AutoExecutor señala que el nivel
adecuado de paralelismo depende de la consulta ejecutada y que una cantidad
inadecuada de executors puede perjudicar el rendimiento o el uso de recursos
\cite{sen2021}. Esto coincide con el comportamiento observado en T1.

Una primera alternativa es utilizar \textbf{pandas}. Para el volumen evaluado,
esta opción obtuvo un tiempo menor y evita la coordinación requerida por un
motor distribuido.

La segunda alternativa consiste en utilizar \textbf{PostgreSQL} para
operaciones de filtrado y agregación cuando el volumen pueda procesarse de
forma adecuada en el servidor de base de datos. Esta opción evita introducir
un componente adicional en la arquitectura para cargas que aún no requieren
procesamiento distribuido.

Spark se reservaría para escenarios donde el crecimiento de los históricos de
pedidos, ventas e inventario supere la capacidad conveniente de una solución
de un solo nodo. Antes de adoptarlo deberían realizarse nuevas pruebas con
diferentes volúmenes para comprobar que el paralelismo compensa su sobrecarga.

Por tanto, Spark puede considerarse una alternativa para el componente
analítico futuro de Tienda Tech, pero las mediciones actuales no justifican
utilizarlo para T1 bajo las condiciones evaluadas.

% Conclusiones
\section{Conclusiones}

La transformación T1 no obtuvo una aceleración al incrementar el número de
executors. Tomando como referencia Spark con un executor, el speedup fue
\num{0.3635} con dos executors y \num{0.9069} con cuatro. Ambos valores fueron
inferiores a uno, por lo que agregar recursos no redujo el tiempo respecto de
la configuración base.

La métrica de Karp-Flatt produjo valores de \num{4.5026} para dos executors y
\num{1.1369} para cuatro. Estos resultados, junto con la dispersión de los
tiempos observada con dos executors, indican que la ejecución estuvo fuertemente
afectada por costos adicionales y que la pérdida de rendimiento no puede
explicarse únicamente mediante una fracción serial constante.

Finalmente, Spark tendría mayor sentido en Tienda Tech como componente de
procesamiento por lotes para históricos de pedidos, ventas e inventario que
como parte del flujo transaccional. Sin embargo, antes de adoptarlo debe
determinarse experimentalmente el volumen a partir del cual la distribución
resulta rentable mediante nuevas mediciones con distintos tamaños de entrada.

% Uso de IA
\section{Declaración de uso de inteligencia artificial generativa}

Se utilizó ChatGPT como herramienta de apoyo para organizar el contenido del
informe de acuerdo con la rúbrica, revisar la redacción técnica, comprobar
operaciones matemáticas y apoyar la estructuración de tablas y figuras.

Los datos experimentales no fueron generados mediante inteligencia artificial.
Los tiempos utilizados proceden de las ejecuciones registradas en el
repositorio PE-U4 del Equipo C. Las mediciones adicionales de T1 con uno y dos
executors se encuentran identificadas en el commit declarado en este informe.

Las conclusiones fueron elaboradas a partir de los resultados experimentales
registrados y las referencias académicas utilizadas se incorporaron para
respaldar únicamente las afirmaciones relacionadas con sus respectivos
resultados y ámbitos de estudio.

% Referencias
\printbibliography

\end{document}